\documentclass{article}
\usepackage{graphicx}
\usepackage{amsmath}
\usepackage{amsfonts}
\usepackage{geometry}
\usepackage{tikz}
\usepackage{syntax}
\usepackage{underscore}
\usepackage{tikz-uml}

% Additional required TikZ libraries
\usetikzlibrary{positioning,shapes,arrows,fit,backgrounds}

% Set margins to accommodate the UML diagram
\geometry{a4paper, margin=1in}

\title{Vehicle Rental Management System Report}
\author{\fontsize{18}{16}\selectfont Latifa Sultan \\[0.5em] \fontsize{12}{14}\selectfont Student ID: 202318117}
\date{\today}

\begin{document}

\maketitle

\section{Phase I: Analysis Report}
\subsection{Application Domain}
The Vehicle Rental Management System is a software application designed to streamline the operations of a vehicle rental business. This system enables users to efficiently manage vehicles, customers, and rental transactions. By automating various processes, the system aims to reduce manual work, minimize errors, and enhance the overall customer experience in vehicle rental services.

\subsection{Functionalities}
The key functionalities of the Vehicle Rental Management System include:

\begin{itemize}
    \item \textbf{Customer Management:} 
    This feature handles customer registration, allowing users to store customer details, track their rental history, and maintain their preferences for future rentals.
    
    \item \textbf{Reservation System:} 
    Users can reserve vehicles in advance, and the system helps avoid double bookings by managing availability effectively.
    
    \item \textbf{Rental Transactions:} 
    This functionality records all aspects of the rental process, including the duration of the rental, extra fees (e.g., for damages), and whether a customer has purchased renter's insurance.
    
    \item \textbf{Payment Processing:} 
    The system can log various payment methods, ensuring a smooth transaction process for customers.
\end{itemize}

\subsection{Initial UML Diagram}
The proposed UML diagram will include the following classes and their relationships:

\begin{itemize}
    \item \textbf{Customer:} Represents the customers of the rental service, containing attributes such as name, contact information, and rental history.
    \item \textbf{Vehicle:} Represents the vehicles available for rent, including attributes like vehicle ID, model, availability status, and rental price.
    \item \textbf{Reservation:} Manages the reservation details, linking customers to the vehicles they wish to rent, including reservation dates and status.
    \item \textbf{RentalTransaction:} Records the details of each rental transaction, including the customer, vehicle, rental duration, and any additional fees.
    \item \textbf{Payment:} Handles payment information, including payment methods and transaction status.
\end{itemize}

\newpage

\section{Phase II: Detailed Report}
\subsection{Project Description}
In this phase, we identify the main classes, objects, methods, and constructors needed for the Vehicle Rental Management System. The following classes are proposed:

\subsection{Main Classes}
\begin{itemize}
    \item \textbf{Customer}
        \begin{itemize}
            \item Attributes: name, age, purchasedRentersInsurance, vehicle, duration, damageFees, paymentMethod, phoneNumber
            \item Methods: 
            \begin{itemize}
                \item +Customer()
                \item +Customer(name, age, insurance, vehicle, duration, phoneNumber)
                \item +setName(String name): void
                \item +getName(): String
                \item +setAge(int age): void
                \item +getAge(): int
                \item +setVehicle(String vehicle): void
                \item +getVehicle(): String
                \item +setDuration(int duration): void
                \item +getDuration(): int
                \item +setReservation(boolean reservation): void
                \item +getReservation(): boolean
                \item +setDamageFees(double n): void
                \item +getDamageFees(): double
                \item +setPaymentMethod(String payment): void
                \item +getPaymentMethod(): String
                \item +setPhoneNumber(String number): void
                \item +getPhoneNumber(): String
            \end{itemize}
        \end{itemize}
    
    \item \textbf{Vehicle}
        \begin{itemize}
            \item Attributes: vehicleID, model, availabilityStatus, rentalPrice
            \item Methods: 
            \begin{itemize}
                \item +Vehicle()
                \item +checkAvailability(): boolean
                \item +updateStatus(): void
                \item +getRentalPrice(): double
            \end{itemize}
        \end{itemize}
    
    \item \textbf{Reservation}
        \begin{itemize}
            \item Attributes: reservationID, customer, vehicle, reservationDate, status
            \item Methods: 
            \begin{itemize}
                \item +Reservation()
                \item +createReservation(): void
                \item +cancelReservation(): void
                \item +updateReservation(): void
            \end{itemize}
        \end{itemize}
    
    \item \textbf{RentalTransaction}
        \begin{itemize}
            \item Attributes: transactionID, customer, vehicle, rentalDuration, extraFees
            \item Methods: 
            \begin{itemize}
                \item +RentalTransaction()
                \item +processTransaction(): void
                \item +calculateTotal(): double
                \item +generateReceipt(): void
            \end{itemize}
        \end{itemize}
    
    \item \textbf{Payment}
        \begin{itemize}
            \item Attributes: paymentID, transaction, paymentMethod, transactionStatus
            \item Methods: 
            \begin{itemize}
                \item +Payment()
                \item +processPayment(): void
                \item +refundPayment(): void
                \item +getPaymentDetails(): string
            \end{itemize}
        \end{itemize}
\end{itemize}

\newpage

\subsection{UML Diagram with Arrows}
Below is the updated UML diagram with directional arrows representing the relationships between classes. 

\begin{center}
\begin{tikzpicture}[scale=0.9, transform shape, node distance=2cm]
    % Customer Class
    \begin{umlclass}[x=0,y=2.5]{Customer}{
        - name: string \\ 
        - age: int \\ 
        - purchasedRentersInsurance: boolean \\ 
        - vehicle: string \\ 
        - duration: int \\ 
        - damageFees: double \\ 
        - paymentMethod: string \\ 
        - phoneNumber: string
    }{
        + Customer() \\ 
        + Customer(name, age, insurance, vehicle, duration, phoneNumber) \\ 
        + setName(name: string): void \\ 
        + getName(): string \\ 
        + setAge(age: int): void \\ 
        + getAge(): int \\ 
        + setVehicle(vehicle: string): void \\ 
        + getVehicle(): string \\ 
        + setDuration(duration: int): void \\ 
        + getDuration(): int \\ 
        + setReservation(reservation: boolean): void \\ 
        + getReservation(): boolean \\ 
        + setDamageFees(fees: double): void \\ 
        + getDamageFees(): double \\ 
        + setPaymentMethod(payment: string): void \\ 
        + getPaymentMethod(): string \\ 
        + setPhoneNumber(number: string): void \\ 
        + getPhoneNumber(): string
    }
    \end{umlclass}

    % Vehicle Class
    \begin{umlclass}[x=-2.7,y=-7.5]{Vehicle}{
        - vehicleID: string \\ 
        - model: string \\ 
        - availabilityStatus: boolean \\ 
        - rentalPrice: double
    }{
        + Vehicle() \\ 
        + checkAvailability(): boolean \\ 
        + updateStatus(): void \\ 
        + getRentalPrice(): double
    }
    \end{umlclass}

    % Reservation Class
    \begin{umlclass}[x=10,y=8]{Reservation}{
        - reservationID: string \\ 
        - customer: Customer \\ 
        - vehicle: Vehicle \\ 
        - reservationDate: string \\ 
        - status: string
    }{
        + Reservation() \\ 
        + createReservation(): void \\ 
        + cancelReservation(): void \\ 
        + updateReservation(): void
    }
    \end{umlclass}

    % RentalTransaction Class
    \begin{umlclass}[x=10,y=0]{RentalTransaction}{
        - transactionID: string \\ 
        - customer: Customer \\ 
        - vehicle: Vehicle \\ 
        - rentalDuration: int \\ 
        - extraFees: double
    }{
        + RentalTransaction() \\ 
        + processTransaction(): void \\ 
        + calculateTotal(): double \\ 
        + generateReceipt(): void
    }
    \end{umlclass}

    % Payment Class
    \begin{umlclass}[x=10,y=-7.5]{Payment}{
        - paymentID: string \\ 
        - transaction: RentalTransaction \\ 
        - paymentMethod: string \\ 
        - transactionStatus: string
    }{
        + Payment() \\ 
        + processPayment(): void \\ 
        + refundPayment(): void \\ 
        + getPaymentDetails(): string
    }
    \end{umlclass}

    % Relationships with arrows
    \umluniassoc[arg1=customer, pos1=0.9, anchor1=west, ->]{Reservation}{Customer}
    \umluniassoc[arg1=vehicle, pos1=0.2, anchor1=west, ->]{Reservation}{Vehicle}
    \umluniassoc[arg1=customer, pos1=0.2, anchor1=west, ->]{RentalTransaction}{Customer}
    \umluniassoc[arg1=vehicle, pos1=0.7, anchor1=west, ->]{RentalTransaction}{Vehicle}
    \umluniassoc[arg1=transaction, pos1=0.7, anchor1=west, ->]{Payment}{RentalTransaction}

    %  % Arrows for associations
    % \umlunicompo{Customer}{Reservation}
    % \umlunicompo{Vehicle}{Reservation}
    % \umlunicompo{Customer}{RentalTransaction}
    % \umlunicompo{Vehicle}{RentalTransaction}
    % \umlunicompo{RentalTransaction}{Payment}

    % % Relationships
    % \umlassoc{Customer}{Reservation}
    % \umlassoc{Vehicle}{Reservation}
    % \umlassoc{Customer}{RentalTransaction}
    % \umlassoc{Vehicle}{RentalTransaction}
    % \umlassoc{RentalTransaction}{Payment}

    

\end{tikzpicture}
\end{center}

\subsection{Relationship Descriptions}
    \begin{enumerate}
        \item \textbf{Customer - Reservation:} The `Reservation` class depends on the `Customer` class to access details about the customer making the reservation. This dependency is represented by an arrow from `Reservation` to `Customer`.
    
        \item \textbf{Vehicle - Reservation: } The `Reservation` class accesses details about the `Vehicle` being reserved. The arrow from `Reservation` to `Vehicle` indicates that `Reservation` holds information about the specific vehicle involved.
    
        \item \textbf{Customer - RentalTransaction: } `RentalTransaction` requires customer details to log transactions accurately. The arrow from `RentalTransaction` to `Customer` shows that `RentalTransaction` depends on `Customer`.
    
        \item \textbf{Vehicle - RentalTransaction: } Since `RentalTransaction` involves a specific vehicle, an arrow from `RentalTransaction` to `Vehicle` indicates this dependency.
    
        \item \textbf{RentalTransaction - Payment: } The `Payment` class needs transaction details (like total amount and transaction ID) from `RentalTransaction` for payment processing. The arrow from `Payment` to `RentalTransaction` illustrates this relationship.
    \end{enumerate}
\end{document}
